\documentclass[12pt,english]{article}
\usepackage[T1]{fontenc}
\usepackage[latin9]{inputenc}
\usepackage[toc,page]{appendix}
\usepackage{geometry}
\geometry{verbose,tmargin=1.5in,bmargin=1.5in,lmargin=1.5in,rmargin=1.5in}
\usepackage{bm}
\usepackage{amstext}
\usepackage{comment}
\usepackage{amsmath}
\usepackage{hyperref}
\hypersetup{
     colorlinks = true,
     citecolor = navy,
     linkcolor = navy,
     urlcolor=navy
}
\usepackage{adjustbox}
\usepackage[space]{grffile}
\usepackage{setspace}
\usepackage[normalem]{ulem}
\usepackage{nicematrix}
\usepackage{xcolor}
\usepackage{soul}
\usepackage{arydshln}
\usepackage{dsfont}
\usepackage{refstyle}
\usepackage[capitalize,nameinlink]{cleveref}
\usepackage{setspace}
\usepackage{bbm}
\usepackage[tableposition=t]{caption} % ensures caption of longtables is normalsized.
\usepackage{mathpazo} % add Palatino font
\usepackage{longtable}
\usepackage{amssymb}
\usepackage{amsfonts}
\usepackage{amsmath}
%\usepackage[margin=1in]{geometry} % to set margins
\usepackage{scalefnt}
\usepackage{wrapfig}
\usepackage{graphicx}
\usepackage{graphics}

%\usepackage{subfigure}
\usepackage{lscape}
\usepackage{pdflscape}
\usepackage{multirow}
\usepackage{array}
\usepackage{rotating}
\usepackage{bibunits}
\usepackage{fix-cm}
\usepackage{float}
\floatplacement{table}{H}
\usepackage{epsfig}
\usepackage{epstopdf}
\usepackage{refstyle}
\usepackage{tikz}
\usepackage{booktabs}
\usepackage{tabularx}
\usepackage{array}
\newcolumntype{x}[1]{>{\raggedleft\arraybackslash}p{#1}}%
\newcolumntype{d}[1]{>{\centering\arraybackslash}p{#1}}%
\newcolumntype{y}[1]{>{\raggedright\arraybackslash}p{#1}}%
\usepackage{verbatim}
%\epstopdfsetup{suffix=}
\setcounter{MaxMatrixCols}{10}
\usepackage[longnamesfirst,authoryear]{natbib}
\newenvironment{proof}[1][Proof]{\noindent \textbf{#1.}}{\rule{0.5em}{0.5em}}
\setlength{\topmargin}{-0.80in} \setlength{\textheight}{9.50in}
\setlength{\evensidemargin}{.0in} \setlength{\oddsidemargin}{.0in}
\setlength{\textwidth}{6.85in}
\setlength{\voffset}{0.4cm}
\setlength{\hoffset}{-0.5cm}
\usepackage{setspace}
\usepackage{afterpage}
\usepackage[flushleft]{threeparttable}
\usepackage{tabulary}
%\usepackage[titletoc,title]{appendix}
\usepackage[bottom]{footmisc}
\usepackage{color}
	\definecolor{maroon}{RGB}{128,0,0}
	\definecolor{navy}{RGB}{0,0,165}
	\definecolor{lb}{RGB}{150,165,190}
	\definecolor{bluenice}{rgb}{0.25,0.51,0.76}
	\definecolor{rednice}{rgb}{0.85,0.26,0}
	\definecolor{orangenice}{HTML}{EB811B}
	\definecolor{greennice}{HTML}{14B03D}
\usepackage{colortbl}
\usepackage{caption}
\usepackage{subcaption}
\usepackage{placeins}
\usepackage{nicefrac}
%\usepackage{chngcntr}
\newtheorem{assumption}{Assumption}
\newtheorem{definition}{Definition}
\usetikzlibrary{matrix,decorations.pathreplacing,arrows.meta,positioning}

\newcommand{\vy}{\underline{y}}
\newcommand{\ve}{\underline{\epsilon}}
\newcommand{\A}[1]{\mathbold{A_{#1}}}
\newcommand{\D}[1]{\mathbold{D_{#1}}}
\newcommand{\C}[1]{\mathbold{C_{#1}}}
\newcommand{\B}[1]{\mathbold{B_{#1}}}
\newcommand{\R}[1]{\mathbold{R_{#1}}}

%\newtheorem{definition}{Definition}
\begin{document}
\title{Consistent Inflation Definitions}
\author{N}
\maketitle

With the Leontief Production we only get the exact match of the two inflations with a very specific definition of inflation.

\[
	\Pi_{CPI,t} \equiv \frac{P_{CPI,t}}{P_{CPI,t-1}} = \frac{\sum_p \omega_{PCE,p,t} P_{p,t}}{\sum_{p'} \omega_{PCE,p',t-1} P_{p',t-1}}
\]

The Shapiro definition of inflation does {\bf NOT} hold exactly any more (because we aren't doing log-approximations).

The goal of this note is to outline the new way to calculate inflation for the graphs. Note that this new way will result in \emph{different} inflation profiles. However, if done for both the Shapiro and ShapirIO graphs, we should still get them summing to the same value.

Let's proceed:

\section{Calculating Inflation}

Note that if we calculate gross inflation the way I said we get the following correspondence:

\begin{align*}
	\Pi_{CPI,t}&\equiv \frac{P_{CPI,t}}{P_{CPI,t-1}} = \frac{\textcolor{orangenice}{\underline{\omega}_{VA,t}'}\underline{P}_{VA,t}}{\textcolor{orangenice}{\underline{\omega}_{VA,t-1}}'\underline{P}_{VA,t-1}} = \frac{\sum_p \textcolor{orangenice}{\omega_{VA,p,t}} P_{VA,p,t}}{\sum_{p'} \textcolor{orangenice}{\omega_{VA,p',t-1}} P_{VA,p',t-1}}& \substack{\text{(Sum notation more}\\\text{convenient for below)}}\\
	&=\sum_p \underbrace{\frac{\textcolor{orangenice}{\omega_{VA,p,t}} P_{VA,p,t-1}}{\sum_{p'} \textcolor{orangenice}{\omega_{VA,p',t-1}} P_{VA,p',t-1}}}_{\equiv \textcolor{bluenice}{\underline{\omega}_{\pi_{VA},p,t}}}\underbrace{\frac{P_{VA,p,t}}{P_{VA,p,t-1}}}_{\equiv \Pi_{VA,p,t}} &\substack{\text{Multiply each term by 1: }\\ \frac{P_{VA,p,t-1}}{P_{VA,p,t-1}}}\\
	&=\sum_{p} \textcolor{bluenice}{\underline{\omega}_{\pi_{VA},p,t}}\Pi_{VA,p,t}
\end{align*}

This introduces an immediate problem: this is \emph{gross} inflation (1.1 = 10\% rise) whereas we want net inflation.

How can we do this with the Shapiro method? See below:

\begin{enumerate}
	\item Calculate Gross inflation for output and value-added prices: $\Pi_{p,t} \equiv P_{p,t}/P_{p,t-1}$ and $\Pi_{VA,p,t}\equiv P_{VA,p,t}/P_{VA,p,t-1}$.
	\item Collect the weights for value added $\omega_{\Pi_{VA},p,t}$ and calculate the following weights for Shapiro:
	\[
		\omega_{\Pi,p,t}\equiv \frac{\omega_{PCE,p,t}P_{p,t-1}}{\sum_{p'}\omega_{PCE,p',t-1}P_{p',t-1}}
	\]
	\item Verify that $\Pi_{CPI,t} = \sum_p \omega_{\Pi,p,t}\Pi_{p,t} = \sum_p\textcolor{bluenice}{\underline{\omega}_{\pi_{VA},p,t}} \Pi_{VA,p,t}$
	\item How then do we get to net inflation? We need the following definitions:
	\begin{itemize}
		\item $\omega_{\Pi,p,t}^{lag}\equiv \frac{\omega_{PCE,p,t-1}P_{p,t-1}}{\sum_{p'}\omega_{PCE,p',t-1}P_{p',t-1}}$
		\item $\textcolor{rednice}{\omega_{VA,p,t}^{lag}}\equiv \frac{\textcolor{orangenice}{\omega_{VA,p,t-1}}P_{VA,p,t-1}}{\sum_{p'}\textcolor{orangenice}{\omega_{VA,p',t-1}}P_{VA,p',t-1}}$
	\end{itemize}	
	
	Calculate supply and demand inflation as:
	\begin{align*}
		\pi_{CPI,t}^{Demand, Shapiro} &= \sum_p D_{p,t} \left[\omega_{\Pi,p,t}\underbrace{\pi_{p,t}}_{\text{NET Inflation}} +  (\omega_{\Pi,p,t} - \omega_{\Pi,p,t}^{lag})\right] \\
				\pi_{CPI,t}^{Supply, Shapiro} &= \sum_p (1-D_{p,t}) \left[\omega_{\Pi,p,t}\underbrace{\pi_{p,t}}_{\text{NET Inflation}} +  (\omega_{\Pi,p,t} - \omega_{\Pi,p,t}^{lag})\right] \\
		\pi_{CPI,t}^{Demand, ShapirIO} &= \sum_p D_{p,t} \left[\textcolor{bluenice}{\omega_{\Pi_{VA},p,t}}\underbrace{\pi_{VA,p,t}}_{\text{NET Inflation}} + (\textcolor{bluenice}{\omega_{\Pi_{VA},p,t}}- \textcolor{rednice}{\omega_{VA,p,t}^{lag}})\right] \\
				\pi_{CPI,t}^{Supply, ShapirIO} &= \sum_p (1-D_{p,t}) \left[\textcolor{bluenice}{\omega_{\Pi_{VA},p,t}}\underbrace{\pi_{VA,p,t}}_{\text{NET Inflation}} + (\textcolor{bluenice}{\omega_{\Pi_{VA},p,t}}- \textcolor{rednice}{\omega_{VA,p,t}^{lag}})\right] \\
	\end{align*}
	\item Verify that 
	\begin{align*}
		\pi_{CPI,t}^{Demand, Shapiro} + \pi_{CPI,t}^{Supply, Shapiro} &=  \pi_{CPI,t} \equiv \Pi_{CPI,t} -1\\
		\pi_{CPI,t}^{Demand, ShapirIO} + \pi_{CPI,t}^{Supply, ShapirIO} &=\pi_{CPI,t} \equiv \Pi_{CPI,t} -1
		\end{align*}
	\item Only then do the 12-month lag thing.
\end{enumerate}


\end{document}